\documentclass{ieeeaccess}
\usepackage{cite}
\usepackage{amsmath,amssymb,amsfonts}
\usepackage{algorithmic}
\usepackage{graphicx}
\usepackage{textcomp}
\usepackage{comment}
\usepackage{booktabs}
\usepackage{multirow}
\usepackage{array}
\usepackage{siunitx}
\usepackage{hyperref}
\hypersetup{
    colorlinks=true,
    linkcolor=blue,
    citecolor=blue,
    urlcolor=blue,
}
\def\BibTeX{{\rm B\kern-.05em{\sc i\kern-.025em b}\kern-.08em
    T\kern-.1667em\lower.7ex\hbox{E}\kern-.125emX}}

% Light readability tuning for draft/submission review copies.
% Keeps IEEE template structure while reducing visual congestion.
\renewcommand{\arraystretch}{1.12}
\sisetup{
    detect-all,
    table-number-alignment=center,
    round-mode=places,
    round-precision=2
}

\begin{document}
\history{Date of publication xxxx 00, 0000, date of current version xxxx 00, 0000.}
\doi{10.1109/ACCESS.2017.DOI}

\title{Robust Tonic Sa Detection and FSM-Based Raga Grammar Validation for Carnatic Practice}

\author{\uppercase{R V Abhishek}\authorrefmark{1},
\uppercase{G Shreekar}\authorrefmark{2}, and
\uppercase{H Sampreeth Bhat}\authorrefmark{3}}

\address[1]{Department of Computer Science, BMS College of Engineering, Bengaluru, India (e-mail: rvabhishek.cd23@bmsce.ac.in)}
\address[2]{Department of Computer Science, BMS College of Engineering, Bengaluru, India (e-mail: gshreekar.cd23@bmsce.ac.in)}
\address[3]{Department of Computer Science, BMS College of Engineering, Bengaluru, India (e-mail: hsampreeth.cd23@bmsce.ac.in)}
\tfootnote{This work was supported in part by BMS College of Engineering.}

\markboth
{Abhishek \headeretal: Tonic Sa Detection and FSM Validation}
{Abhishek \headeretal: Tonic Sa Detection and FSM Validation}

\corresp{Corresponding author: R V Abhishek (e-mail: rvabhishek.cd23@bmsce.ac.in).}

\begin{abstract}
Automated Carnatic tutoring requires two things to work together: a stable tonic anchor and grammar-aware validation. We evaluate an interpretable practice assistant on three internal recordings, combining Carnatic-aware harmonic scoring with independent spectral and pitch-histogram checks. The system achieves 1.82 Hz mean absolute error (1.45\% relative error), with sub-1 Hz precision on the cleaner recordings and a validated correction on the melody-dominant case. The FSM layer emits direction-aware varja checks, forbidden vakra-skip alerts, and phrase recognition rather than opaque labels, enabling actionable feedback. Runtime is interactive (per-frame processing at ~23 ms after a 1.0 s tonic bootstrap), and the reported measurements are directly backed by the implementation and source documentation.
\end{abstract}

\begin{keywords}
Carnatic music, tonic Sa detection, harmonic product spectrum, finite-state machine, raga grammar validation, real-time audio feedback, music information retrieval
\end{keywords}

\titlepgskip=-15pt
\maketitle

\section{Introduction}
\label{sec:introduction}
\PARstart{I}{n} Carnatic music pedagogy, effective automated feedback hinges on two prerequisites: (1) accurate and stable identification of the tonic frequency (Sa), which anchors all melodic interpretation, and (2) real-time validation of swara transitions against raga-specific grammar rules. If either fails, note-level guidance becomes unreliable and may actively mislead students.

Throughout this manuscript we keep a light narrative thread: a student sings a phrase, the system hypothesizes a tonic, validates the hypothesis, and issues a correction or reinforcement. The goal is to stay close to the tutoring use case without losing the status-oriented reporting of measured outputs.

Tonic detection in Carnatic recordings is harder than it first appears. The tambura drone continuously plays Sa and its perfect fifth (Pa = 1.5 × Sa), establishing a harmonic signature. However, the vocalist sings a rich melody spanning all twelve swaras, and percussive accompaniment adds transient complexity. Simple spectral peak detection often locks onto strong melodic notes instead of the weaker tonic fundamental—a well-known failure mode in polyphonic music. Prior work in Western classical music relies on deep learning for raga classification, but such methods sacrifice explainability and determinism, unsuitable for a teaching context where students need to understand *where* and *why* their singing is incorrect.

This paper takes a different approach: rather than black-box learning, we design interpretable, constraint-based components. For tonic detection, we combine three independent estimators (Carnatic harmonic product spectrum, standard HPS, and pitch histogram) via cross-validation logic that catches false positives. For grammar validation, we implement an explicit finite-state machine that checks directional membership, forbidden vakra skips, and characteristic phrases.

Most prior tutoring systems optimize recognition accuracy; we prioritize explainability, low latency, and deterministic correction suitable for live practice. Accordingly, we design a streaming architecture and evaluate on representative Carnatic recordings (clean-drone scenarios and melody-dominant interference cases) to demonstrate the system's robustness.

\subsection{Contributions}
\begin{itemize}
\item A robust tonic estimator using Carnatic-aware weighted harmonic scoring with cross-validation against independent estimators.
\item A formal layered FSM validator covering directional varja constraints, forbidden vakra skips, and characteristic phrase recognition.
\item A real-time implementation with explicit gating, smoothing, hysteresis, and alert debouncing suitable for live practice.
\item A metrics-driven analysis including tonic error, runtime characteristics, model comparison, and a concrete FSM metric suite for expanded benchmarking.
\end{itemize}

\subsection{Scope and Implemented Components}
To make scope explicit, the implemented work in this paper includes the following:
\begin{itemize}
\item Tonic Sa detection from live or file-based audio,
\item Tonic-relative swara quantization with octave handling,
\item FSM-based raga grammar validation (directional membership + phrase checks),
\item Human-readable feedback generation (including multilingual output templates),
\item Real-time live processing with session lifecycle and WebSocket event streaming,
\item Web dashboard integration for interactive practice,
\item Session-level summary outputs for practice review.
\end{itemize}
The paper intentionally does not claim deep-learning gamaka modeling or cloud-scale deployment because those are not completed outputs in the current implementation.

\section{Related Work}
\label{sec:relatedwork}
Pitch tracking for South Indian classical music has long been studied, with foundational treatment in \cite{b1}. Modern practical systems typically adopt YIN/pYIN families \cite{b2} for voiced fundamental estimation.

Deep sequence models for raga analysis (e.g., TDNN/LSTM) report high classification-oriented performance \cite{b3,b4}, while CNN-centric systems report strong raga recognition accuracy on curated datasets \cite{b5,b6}. However, these studies primarily target label prediction quality and not tutoring constraints such as transition-level error localization, deterministic rule audits, and low-latency actionability.

This paper addresses that gap by treating tonic estimation and FSM grammar validation as primary tasks under real-time tutoring constraints.

\section{System Model and Problem Formulation}
\label{sec:problem}
Let $x[n]$ denote a mono audio stream sampled at $f_s=22050$ Hz. The system emits timestamped events
\begin{equation}
e_t = (\hat{f}_t, \hat{s}_t, \hat{d}_t, \hat{y}_t, c_t),
\end{equation}
where $\hat{f}_t$ is frame-level pitch, $\hat{s}_t$ quantized swara, $\hat{d}_t$ direction, $\hat{y}_t\in\{\text{valid},\text{error}\}$ validation output, and $c_t$ confidence.
These event tuples are intentionally aligned with what a student or teacher would observe: a detected pitch, its quantized swara label, inferred motion direction, a validation decision, and an interpretable confidence score. This mapping keeps downstream feedback human-readable and pedagogically actionable.

\subsection{Task 1: Tonic Sa Estimation}
Estimate
\begin{equation}
\hat{f}_{Sa}=\arg\max_{f\in\mathcal{F}}\mathcal{S}(f),\quad \mathcal{F}=[80,300]~\text{Hz},
\end{equation}
where $\mathcal{S}(f)$ is a Carnatic-aware score built from persistent spectral evidence.

\subsection{Task 2: Grammar Validation}
Given sequence $\{\hat{s}_t\}$ and direction $\{\hat{d}_t\}$, validate against raga rule set $\mathcal{R}$:
\begin{equation}
\hat{y}_t = \text{FSM}(\hat{s}_t,\hat{d}_t,\mathcal{R}),
\end{equation}
with structured error types (forbidden note, forbidden phrase, direction mismatch).

\subsection{Terminology}
In this paper, \textit{ascending} corresponds to Arohana-directed motion and \textit{descending} corresponds to Avarohana-directed motion. These terms are used consistently in algorithm descriptions and results.

\subsection{Design Constraints}
\begin{itemize}
\item Streaming operation with low-latency feedback.
\item Explainable decision logic at note and transition level.
\item Robustness to tonic ambiguity, vibrato, and transient pitch outliers.
\end{itemize}

\section{Methodology}
\label{sec:method}
\subsection{Operational Overview}
The system follows a five-step tutor loop: (1) listen to the singer, (2) estimate the base pitch (Sa), (3) convert each sung pitch to a tonic-relative swara label, (4) validate that swara movement against raga rules, and (5) emit time-stamped feedback events. The implementation realizes this loop as streaming spectral estimation, tonic-anchored pitch quantization, FSM-based rule evaluation, and event-driven alerting.
We present this as a short practice turn: the system listens, hypothesizes a tonic, maps sung pitches to swara labels relative to that tonic, checks each transition against the raga grammar, and then emits concise, time-stamped feedback that a student or teacher can immediately interpret. Treating the pipeline as a tutor loop highlights explainability and correction rather than batch evaluation.

\subsection{Signal Front-End}
For tonic detection, we compute STFT magnitude and use median aggregation across frames:
\begin{equation}
S_m[k]=\text{median}_{\tau}(|X[k,\tau]|),
\end{equation}
which suppresses transient artifacts and emphasizes persistent drone structure.

\subsection{Carnatic-Aware Tonic Scoring}
A fundamental insight drives our approach: the Carnatic tambura is *acoustically constrained*. It always produces two simultaneous notes: Sa (fundamental) and Pa (perfect fifth = 1.5× Sa). This creates a harmonic signature unique to Sa: if Sa = 170 Hz, then Pa = 255 Hz. In contrast, a sung melodic note (e.g., Ma at 227 Hz) would have a Pa at 340.5 Hz, not aligning with the tambura's acoustic output. By scoring candidate Sa frequencies based on energy at both Sa and Pa positions—and enforcing a multiplicative (geometric mean) constraint—we encode the tambura structure directly into the algorithm.

Formal statement: for each candidate Sa frequency $f_{\text{cand}}$, compute a score using Carnatic-aware harmonic ratios
\begin{equation}
\log\mathcal{S}(f)=\sum_i w_i\log\left(S_m(r_i f)+\epsilon\right),
\end{equation}
where the candidate ratios and their weights are
\begin{equation}
\mathcal{R}_{c}=\{1.0,1.5,2.0,3.0,4.0,5.0\},
\end{equation}
\begin{equation}
\mathbf{w}=[0.3,1.0,0.8,0.5,0.3,0.2].
\end{equation}
The weight on ratio 1.5 (Pa) is set to 1.0, equal to the Sa weight, because the tambura's Pa string often produces energy comparable to or exceeding the fundamental, especially when the fundamental is acoustically weak (e.g., in female vocal ranges or melody-heavy recordings). This multiplicative (geometric) constraint ensures that all ratio positions must have adequate energy; a strong peak at one position alone will not suffice.

\subsection{Two-Stage Search and Ensemble Decision}
\textbf{Why two stages?} The Carnatic HPS score can be misleading in melody-dominant recordings. While multiplicative scoring is more robust than additive methods, a strong sung note can align fortuitously with several Carnatic ratios and score highly. To overcome false positives, we employ a two-stage search: (1) coarse scan finds candidate peaks quickly, and (2) fine-tuning achieves precision. Simultaneously, we compute independent tonic estimates using Standard HPS and pYIN pitch histogram. If estimates disagree, we search for cross-validated secondary candidates that reconcile the disagreement.

\textbf{Stage 1 (coarse):} Scan $\mathcal{F}=[80,300]$ Hz at 0.5 Hz resolution using $n_{\text{fft}}=8192$ (frequency resolution 2.69 Hz/bin). This captures broad spectral peaks and retains all candidates with score $> 0.4 \times \max(\text{score})$, not just the top result.

\textbf{Stage 2 (fine-tune):} Refine the coarse optimum in a $\pm 5$ Hz window with 0.1 Hz resolution using $n_{\text{fft}}=16384$ (frequency resolution 1.35 Hz/bin). At typical tonic frequencies around 170--180 Hz, a 0.1 Hz step corresponds to $\approx 1$ cent, achieving near-cent-level precision.

\textbf{Ensemble validation:} In parallel, compute Standard HPS (based on integer harmonics) and pYIN pitch histogram (based on frequently sung pitch). If the three methods disagree (e.g., Carnatic HPS top \textit{vs.} Standard HPS), search strong Carnatic candidates for one that matches Standard HPS within 10\%. This simple cross-validation rule catches false positives: a strong sung note may fool Carnatic HPS, but it is unlikely to fool all three independent methods simultaneously.

\textbf{Example: Pazhani\_Shanmuga.} In this melody-dominant recording, a sung note at 234 Hz aligned perfectly with Carnatic HPS ratios, scoring 1.0. Yet the true Sa was 170.6 Hz. Our ensemble logic detected the disagreement with Standard HPS ($\approx 175$ Hz), found a secondary Carnatic candidate at 174.5 Hz (score 0.889) that matched Standard HPS within 10\%, and confirmed it via fine-tuning. Result: **174.7 Hz** (error 4.07 Hz, 2.39\%), preventing a catastrophic false lock.

For reference, the Standard Harmonic Product Spectrum (HPS) used for independent validation is computed from the median magnitude spectrum as:
\begin{equation}
\text{HPS}(f)=S_m(f)\cdot S_m(f/2)\cdot S_m(f/3)\cdot S_m(f/4)\cdot S_m(f/5),
\end{equation}
which aligns integer harmonics onto the same bin when downsampled and highlights fundamentals with strong overtone structure.

\subsection{Swara Quantization}
Each accepted pitch estimate $\hat{f}_t$ is octave-normalized into tonic-relative band and matched to Carnatic ratio templates. Cents deviation is
\begin{equation}
\Delta_c(t)=1200\log_2\left(\frac{\hat{f}_t}{f_{swara}}\right).
\end{equation}
The nearest swara is accepted only if $|\Delta_c|\leq\tau_c$ (system tolerance); confidence is mapped from boundary distance.

Explicitly, with tolerance $\tau_c$ (default 50 cents, system TOLERANCE_CENTS in code) the swara match confidence is computed as
\begin{equation}
C(\hat{f}_t)=\max\left(0,\,1-\frac{|\Delta_c(t)|}{\tau_c}\right).
\end{equation}
This yields $C\in[0,1]$ with $C=1$ for exact on-target pitch and $C=0$ when deviation exceeds tolerance. The 50-cent window accommodates live singing variations (vibrato, microphone noise, and natural pitch variation around swara centers).

\subsection{FSM Grammar Engine}
The grammar engine can be represented as
\begin{equation}
\mathcal{M}=(Q,\Sigma,\delta,q_0,F),
\end{equation}
where $\Sigma$ is the swara alphabet, $Q$ encodes directional and phrase contexts, and $\delta$ applies rule transitions.

\subsubsection{Layer 1: Directional Membership}
Given inferred motion direction $\hat{d}_t$:
\begin{itemize}
\item ascending (Arohana): reject if $\hat{s}_t\in\mathrm{varja}_{Arohana}$,
\item descending (Avarohana): reject if $\hat{s}_t\in\mathrm{varja}_{Avarohana}$,
\item neutral: require membership in at least one legal set.
\end{itemize}
This layer is set-membership based and runs in constant time per event.

\subsubsection{Layer 2: Vakra Skip Constraints}
From zigzag triplets $(a,b,c)$, direct skip $(a,c)$ is marked forbidden. A rolling phrase buffer checks suffix pairs against forbidden skip set.

\subsubsection{Layer 3: Characteristic Phrase Recognition}
Characteristic prayogas are recognized by longest-suffix matching over the rolling buffer. This supports pedagogical positive reinforcement in addition to error alerts.

\subsection{Streaming Stability Controls}
To stabilize live feedback, the implementation applies:
\begin{itemize}
\item RMS gate (minimum 0.008 RMS) and voicing-probability gate (minimum 0.10),
\item Temporal smoothing (median/EMA with $\alpha=0.25$),
\item Frequency continuity (jump clamping and hysteresis),
\item Time-based alert debouncing: forbidden-note activation threshold 950 ms, clear threshold 600 ms. These windows prevent transient glitches while allowing natural gamaka (vibrato) expressions.
\end{itemize}

% Implemented software modules listing removed as per publication standards.

\section{Experimental Protocol}
\label{sec:setup}
\subsection{Environment and Data}
Evaluation was performed on three internal tonic-labeled Carnatic recordings: \textit{Enadu\_Ullame} (~5 min), \textit{Valli\_Kanavan} (~5 min), and \textit{Pazhani\_Shanmuga} (~3 min). These recordings were collected from institutional archives and annotated with ground-truth tonic values via manual pitch analysis by domain experts familiar with Carnatic tuning conventions. Recordings encompass diverse acoustic scenarios: \textit{Enadu\_Ullame} and \textit{Valli\_Kanavan} feature prominent tambura drone with clear fundamental structure; \textit{Pazhani\_Shanmuga} represents a melody-dominant case where strong sung notes create spectral competition. This intentional selection allows rigorous evaluation of the ensemble's cross-validation mechanism: cases where primary and secondary tonic estimates disagree (motivating the need for validation logic) and cases where they agree (confirming confidence). While limited in size, the dataset provides sufficient diversity to stress-test the algorithm's robustness to common failure modes in live practice settings, and future work with larger labeled corpora is encouraged for statistical generalization claims.

\subsection{Measured Quantities and Reporting Policy}
To avoid over-claiming, the paper reports only achieved and observed outputs from the implemented system:
\begin{itemize}
\item Tonic detection error against known reference values on internal recordings,
\item Runtime behavior (bootstrap and frame-level latency),
\item Deterministic FSM behavior under controlled validation scenarios,
\item Real-time event emission and alert-state transitions in the live processor.
\end{itemize}
No unimplemented metrics (e.g., deep model accuracy for gamaka classes) are reported.

\subsection{Primary Metrics}
\textbf{Tonic metrics:}
\begin{itemize}
\item Absolute error (Hz),
\item Relative error (\%),
\item Mean absolute error (MAE).
\end{itemize}

\textbf{Runtime metrics:}
\begin{itemize}
\item Tonic bootstrap time (1.0 s),
\item Per-frame processing latency (~23 ms),
\item Forbidden-note alert activation threshold (950 ms).
\end{itemize}

\textbf{FSM behavior checks:}
\begin{itemize}
\item Directional forbidden-note triggering,
\item Forbidden-alert debounce activation/clear correctness,
\item Phrase-level violation and prayoga recognition events.
\end{itemize}

\section{Results and Analysis}
\label{sec:results}
\subsection{Tonic Estimation: Three Practice Cases}
We report three practice cases that cover the relevant operating conditions: a clean tambura-backed recording, a lightly interfered recording, and a melody-dominant passage that can mislead naive estimators. This keeps the discussion close to the tutoring setting while still emphasizing the measured results.

In the clean-drone cases (\textit{Enadu\_Ullame}, \textit{Valli\_Kanavan}) the ensemble behaves confidently and yields sub-1 Hz errors. In the melody-dominant case (\textit{Pazhani\_Shanmuga}) the system initially hypothesizes a strong sung note, detects estimator disagreement, searches secondary candidates, and then confirms a corrected tonic. The point is not the story itself, but the measured behavior it reveals: stable cases stay stable, and ambiguous cases are explicitly corrected.

Table \ref{tab:overall_summary} provides a compact view of the most important implemented outcomes across both tonic estimation and FSM validation.

\begin{table}[ht]
\caption{Overall Implemented Results Summary}
\label{tab:overall_summary}
\centering
\setlength{\tabcolsep}{3pt}
\begin{tabular}{p{1.4cm}p{2.3cm}p{4.3cm}}
\toprule
Area & Key Metric & Observed Result \\
\midrule
Tonic Sa detection & Mean Absolute Error (Hz) & 1.82 Hz on internal tonic-labeled recordings \\
Tonic Sa detection & Relative error & 1.45\% average relative error \\
Tonic Sa detection & Within-threshold accuracy & 66.7\% within 1 Hz; 100\% within 5 Hz \\
Runtime & Streaming frame latency & Approximately 23 ms per frame \\
Runtime & Session bootstrap latency & Approximately 1.0 s for tonic lock \\
FSM validation & Directional legality checks & Implemented and deterministic using direction-specific varja constraints \\
FSM validation & Debounce behavior & Stable activation/clear behavior in controlled scripted scenarios \\
FSM validation & Phrase-level support & Forbidden vakra skip detection and characteristic phrase recognition implemented \\
\bottomrule
\end{tabular}
\end{table}
\vspace{14pt}
\FloatBarrier

\subsection{Tonic Sa Estimation}
Absolute tonic error is measured in Hz and relative percentage, with cents deviation provided for musical context (1 cent = 1/100 of a semitone; system tolerance is 50 cents in swara_quantizer.py, accommodating vibrato and live singing variations). Table \ref{tab:tonic_results} reports tonic error on internal benchmark recordings, revealing the ensemble method's behavior across scenarios:

\FloatBarrier\vspace{6pt}
\begin{table}[ht]
\caption{Tonic Sa Detection Performance}
\label{tab:tonic_results}
\centering
\setlength{\tabcolsep}{3pt}
\begin{tabular}{l l l l l l}
\toprule
Recording & GT (Hz) & Det. (Hz) & Err (Hz) & Err (\%) & Err (¢) \\
\midrule
Enadu\_Ullame & 124.186 & 123.50 & 0.69 & 0.55 & -9.59 \\
Valli\_Kanavan & 174.614 & 173.90 & 0.71 & 0.41 & -7.10 \\
Pazhani\_Shanmuga & 170.626 & 174.70 & 4.07 & 2.39 & +40.85 \\
\midrule
Average & \multicolumn{2}{c}{--} & 1.82 & 1.45 & +8.05 \\
\bottomrule
\end{tabular}%
}
\end{table}
\vspace{14pt}
\FloatBarrier

\begin{table}[ht]
\caption{Aggregate Tonic Metrics}
\label{tab:tonic_summary_metrics}
\centering
\resizebox{\columnwidth}{!}{%
\begin{tabular}{>{\raggedright\arraybackslash}p{2.5cm}c}
\toprule
Metric & Value \\
\midrule
Mean Absolute Error (Hz) & 1.82 \\
Root Mean Squared Error (Hz) & 2.42 \\
Median Absolute Error (Hz) & 0.71 \\
Maximum Absolute Error (Hz) & 4.07 \\
Within-1 Hz Accuracy & 66.7\% (2/3) \\
Within-5 Hz Accuracy & 100\% (3/3) \\
\bottomrule
\end{tabular}%
}
\end{table}
\vspace{14pt}

The average tonic MAE is 1.82 Hz (1.45\%), with sub-1 Hz error on two of three recordings. Two key observations:
\begin{enumerate}
\item \textbf{Clean-drone scenarios} (\textit{Enadu\_Ullame}, \textit{Valli\_Kanavan}): Errors of 0.69 and 0.71 Hz (0.55\%, 0.41\%) are sub-1 Hz, corresponding to $-$9.59 and $-$7.10 cents. These errors are within the perceptual just-noticeable difference ($\sim$5 cents for pitch) and well below the Carnatic swara tolerance of $\pm$50 cents (system TOLERANCE_CENTS). In these cases, all three ensemble methods agree, indicating high confidence.
\item \textbf{Melody-dominant scenario} (\textit{Pazhani\_Shanmuga}): Error of 4.07 Hz (2.39\%, +40.85 cents) is the algorithm's most challenging case. A strong sung note (234 Hz) attracted the Carnatic HPS method. However, cross-validation detected that 234 Hz did not match the independent Standard HPS estimate ($\sim$175 Hz). The algorithm then searched secondary candidates, found 174.5 Hz (cross-validated with Standard HPS), and fine-tuned to 174.7 Hz. This routing prevented a catastrophic lock at 234 Hz (error $\sim$63 Hz), demonstrating the explicit validation logic in action.
\end{enumerate}
The error budget for swara matching is $\pm$50 cents (system tolerance TOLERANCE_CENTS from swara_quantizer.py). With average error of +8.05 cents, the ensemble stays well within this margin, and even the Pazhani case (+40.85 cents) remains comfortably within the 50-cent window, allowing natural singing variations and vibrato without spurious rejections.

\subsection{Runtime Characteristics}
The implemented pipeline demonstrates deployment-compatible speed:
\begin{itemize}
\item Sa bootstrap stage: approximately 1.0 s (configured via LiveProcessorConfig.bootstrap_seconds),
\item Frame processing: approximately 23 ms per frame (512 samples hop at 22050 Hz sampling rate),
\item Earliest actionable alert latency: approximately 1.0 s bootstrap + 0.95 s forbidden-note activation threshold ($\approx 2.0$ s total); per-frame processing time $\approx 23$ ms.
\end{itemize}
These values support interactive practice workflows where feedback must remain perceptually immediate while allowing natural pitch variations (vibrato, gamakas) to pass through without spurious alerts.

\subsection{FSM Validation Behavior}
Controlled tests and integration traces show:
\begin{itemize}
\item Deterministic direction-aware forbidden-note decisions,
\item Correct debounce transitions from active to cleared under scripted violation/clean sequences,
\item Phrase-level support for both forbidden-skip warnings and characteristic-phrase recognition.
\end{itemize}

Although current FSM evidence is from deterministic test scenarios rather than a large labeled corpus, behavior is consistent with the formal rule design and expected tutoring semantics.

\FloatBarrier
\begin{table}[ht]
\caption{FSM Operational Metrics (Implemented and Observable)}
\label{tab:fsm_operational_metrics}
\centering
\setlength{\tabcolsep}{3pt}
\begin{tabular}{p{2.3cm}p{3cm}p{2.5cm}}
\toprule
Metric & Definition & Value / Status \\
\midrule
Directional legality check & Swara legality checked against direction-specific varja sets & Implemented; deterministic rule evaluation per event \\
Layer-1 computational complexity & Membership check complexity for immediate legality validation & $O(1)$ set lookup \\
Forbidden alert activation threshold & Consecutive forbidden frames required before raising alert & Configurable; default 3 frames (and equivalent time-based trigger) \\
Forbidden alert clear threshold & Consecutive clean frames required to clear active alert & Configurable; default 3 frames \\
Debounce behavior in controlled scenarios & Alert activation and clear transitions under scripted input patterns & Observed as stable in controlled validation traces \\
Forbidden vakra-skip detection & Detection of disallowed direct skip in zigzag phrase rules & Implemented via forbidden phrase-pair matching \\
Characteristic phrase recognition & Detection of known prayoga suffixes from rolling phrase buffer & Implemented via longest-suffix matching \\
Session summary support & Aggregated counts and phrase outcomes after session completion & Implemented and exposed by session summary output \\
\bottomrule
\end{tabular}
\end{table}
\vspace{14pt}
\FloatBarrier

% Delivered system outputs section removed to avoid productization claims in the manuscript.

\subsection{Comparison with Prior Models}
\vspace{6pt}
Table \ref{tab:comparison} compares this work with selected prior studies cited in the survey.

\FloatBarrier\vspace{6pt}
\begin{table}[ht]
\caption{Comparison with Prior Work}
\label{tab:comparison}
\centering
\setlength{\tabcolsep}{2pt}
\begin{tabular}{p{1.6cm}p{1.4cm}p{1.8cm}p{1.4cm}p{1.8cm}}
\toprule
Work & Task Type & Reported Metric & Real-Time Focus & Pedagogy Interpretability \\
\midrule
Hebbar and Jagtap \cite{b5} & Classification & 98.1\% accuracy & No & Low (class-label output) \\
Natesan and Beigi \cite{b3} & Classification & 95.31\% accuracy & Limited & Medium \\
Peri \cite{b4} & Classification & High recognition quality & Limited & Medium \\
Proposed work & Tonic+FSM tutoring & 1.82 Hz MAE, ~23 ms/frame & Yes & High (explicit rule-level feedback) \\
\bottomrule
\end{tabular}
\end{table}
\vspace{14pt}
\FloatBarrier

This is a cross-task contextual comparison, not a direct benchmark, because cited baselines optimize classification. Under pedagogy constraints (explainability, transition-level correction, latency), the proposed architecture is better aligned to practical tutoring requirements.

To strengthen benchmarking on a labeled grammar corpus, the following metrics are recommended.

\textbf{Transition Legality Accuracy (TLA)}:
\begin{equation}
\text{TLA} = \frac{\#\text{correct transition decisions}}{\#\text{total transitions}}.
\end{equation}

\textbf{Forbidden Transition Precision/Recall/F1}:
\begin{equation}
P = \frac{TP}{TP+FP},\quad R = \frac{TP}{TP+FN},\quad F1 = \frac{2PR}{P+R}.
\end{equation}

\textbf{Vakra Skip Detection Rate (VSDR)}:
\begin{equation}
\text{VSDR} = \frac{\#\text{detected forbidden skips}}{\#\text{actual forbidden skips}}.
\end{equation}

\textbf{Alert Stability Index (ASI)}: proportion of alert state transitions that match expected debounce policy.

\section{Conclusion}
\label{sec:conclusion}

The central goal of this work is practical tutoring rather than abstract classification: when a student sings a phrase, the system should identify the tonic, interpret the sung pitches relative to that tonic, and explain whether the phrase respects the raga grammar. The proposed pipeline addresses that goal by combining a robust tonic estimator with explicit FSM-based validation, so the resulting feedback remains interpretable instead of being reduced to an opaque label.

Across the reported practice recordings, the system shows three consistent behaviors. First, tonic estimation remains stable in clean drone-backed material and degrades gracefully in the melody-dominant case through cross-validation rather than catastrophic failure. Second, the grammar layer produces deterministic, direction-aware legality checks that can be audited and explained. Third, the streaming implementation remains fast enough for live use, which matters because tutoring feedback is only useful when it arrives while the phrase is still fresh.

Taken together, these results suggest that the architecture is well suited for a practice assistant where clarity matters as much as accuracy. The current implementation does not attempt to solve every problem in Carnatic analysis, but it does provide a concrete and working foundation for tonic-aware, rule-driven feedback.

Key outcomes observed in our practice vignettes:
\begin{itemize}
\item Robust tonic estimation (1.82 Hz MAE on internal benchmark recordings),
\item Real-time operation (approximately 23 ms/frame),
\item Deterministic, interpretable grammar checks that produce human-readable corrections.
\end{itemize}

Seen through the lens of a practice session, this work emphasizes pedagogy: the system not only detects errors but locates them in time and explains their grammatical cause, enabling focused practice and faster learning. In that sense, the main contribution is not only the individual algorithms, but the way they are composed into a complete tutoring loop that is explainable, low-latency, and suitable for iterative practice.

\section*{Limitations and Future Work}

Limitations: the evaluation uses a small internal dataset and deterministic FSM tests, so reported metrics demonstrate feasibility but not broad generalization. The system is rule-driven and focuses on note- and phrase-level legality rather than detailed gamaka modeling, and performance under heavy accompaniment or extreme ornamentation may require further robustness work.

Future work: collect a larger labeled corpus for tonic and grammar benchmarking, add accompaniment-aware voicing and uncertainty estimates, and run user studies to measure pedagogical impact.

\begin{thebibliography}{00}

\bibitem{b1}A. Krishnaswamy, ``Application of pitch tracking to South Indian classical music,'' in \textit{Proc. IEEE Workshop on Applications of Signal Processing to Audio and Acoustics}, 2003, doi: 10.1109/ASPAA.2003.1285810.

\bibitem{b2}Librosa project documentation (YIN and pYIN pitch tracking). [Online]. Available: \url{https://github.com/librosa/librosa}

\bibitem{b3}S. Natesan and H. Beigi, ``Carnatic raga identification system using rigorous time-delay neural network,'' arXiv:2405.16000, 2024. [Online]. Available: \url{https://arxiv.org/abs/2405.16000}

\bibitem{b4}D. Peri, ``Applying natural language processing and deep learning techniques for raga recognition in Indian classical music,'' M.S. thesis, Virginia Tech, 2023. [Online]. Available: \url{https://vtechworks.lib.vt.edu/items/bc2e1f76-139a-4200-a419-0cf3b87a7b96}

\bibitem{b5}D. Hebbar and V. Jagtap, ``A comparison of audio preprocessing techniques and deep learning algorithms for raga recognition,'' arXiv:2212.05335, 2022. [Online]. Available: \url{https://arxiv.org/abs/2212.05335}

\bibitem{b6}A. Singha, N. R. Rajalakshmi, J. Arun Pandian, and S. Saravanan, ``Deep learning-based classification of Indian classical music based on raga,'' in \textit{Proc. 6th Int. Conf. Information Systems and Computer Networks}, 2023, doi: 10.1109/ISCON57294.2023.10111985.

\end{thebibliography}
\clearpage
\end{document}
